\clearpage\thispagestyle{empty}
\pdfbookmark[0]{Quelle zu von Lindenaus Brief vom 17. Februar 1806}{lindenau1806letter-sources}
\begin{center}\textbf{Quelle zu von Lindenaus Brief vom 17. Februar 1806}\end{center}
\noindent Diese KI-Übertragung ins heutige Deutsch folgt dem Auszug auf S.~486 der \href{https://gdz.sub.uni-goettingen.de/id/PPN236007467}{\textit{Werke}, Band~VI} (Göttingen, 1874), die den Brief als Teil von Gauß’ Nachlass ausweist. Bernhard von Lindenau schrieb am 17. Februar 1806 aus Altenburg an Carl Friedrich Gauß.
\par\smallskip\noindent Lindenau spricht darüber, dass Franz Xaver von Zach Gauß’ Brief abgedruckt hat, und erläutert seinen eigenen Vorbehalt gegenüber Rohdes Methode. Die Abkürzung Obh.\ folgt der Werkausgabe. KI-Übertragung und KI-Satz.
\clearpage
